\section{AMC 8 2016}
        \begin{enumerate}
        	\item (\textit{AMC 8 2016}) The longest professional tennis match lasted a total of 11 hours and 5 minutes. How many minutes is that?


 $\textbf{(A) } 605 \qquad\textbf{(B) } 655\qquad\textbf{(C) } 665\qquad\textbf{(D) } 1005\qquad \textbf{(E) } 1105$


 
 \textbf{Answer: }C



\textbf{Solution 1}

It is best to split 11 hours and 5 minutes into 2 parts, one of 11 hours and another of 5 minutes.  We know that there is $60$ minutes in a hour.  Therefore, there are $11 \cdot 60 = 660$ minutes in 11 hours.  Adding the second part(the 5 minutes) we get $660 + 5 = \boxed{\textbf{(C)}\ 665}$.

 



\textbf{Solution 2}

The best method comes when you remember your multiplication tables. Thus trivial, we get our answer $\boxed{\textbf{(C) }665}$.

 



\textbf{Solution 3}

11 hours 5 min = $(11 \cdot 60) + 5 \text{min}  = 665 \text{min}$, therefore $\boxed{\textbf{(C) }665}$.

 Carmen

 


	\item (\textit{AMC 8 2016}) In rectangle $ABCD$, $AB=6$ and $AD=8$.  Point $M$ is the midpoint of $\overline{AD}$.  What is the area of $\triangle AMC$?


 
                    \begin{figure}[H]
                        \centering
                        \includegraphics[width=0.5\textwidth]{Problem/AMC_8_AMC_2016_2_Problem_Figure_1.png}
                    \end{figure}
                    


 $\textbf{(A) }12\qquad\textbf{(B) }15\qquad\textbf{(C) }18\qquad\textbf{(D) }20\qquad \textbf{(E) }24$


 
 \textbf{Answer: }A



\textbf{Solution 1}

Using the triangle area formula for triangles: $A = \frac{bh}{2},$ where $A$ is the area, $b$ is the base, and $h$ is the height. This equation gives us $A = \frac{4 \cdot 6}{2} = \frac{24}{2} =\boxed{\textbf{(A) } 12}$.

 



\textbf{Solution 2}

A triangle with the same height and base as a rectangle is half of the rectangle's area.  This means that a triangle with half of the base of the rectangle and also the same height means its area is one quarter of the rectangle's area.  Therefore, we get $\frac{48}{4} =\boxed{\textbf{(A) } 12}$.

 



\textbf{Solution 3(a check)}

We can find the area of the entire rectangle, DCBA to be $8 \cdot 6=48$ and find DCM area to be $\frac{6 \cdot 4}{2} = 12$ and BCA to be $\frac{6 \cdot 8}{2}=24$, $48-12-24=$ $\boxed{\textbf{(A) } 12}$.

 



\textbf{Solution 4}

A triangle is half of a rectangle. So since M is the midpoint of DA, we can see that triangle DAC is half of the whole rectangle. And it is also true that triangle MAC is half of triangle DAC, so triangle MAC would just be 1/4 of the whole rectangle. Since the rectangle's area is 48, 1/4 of 48 would be 12. Which gives us the answer $\boxed{\textbf{(A) } 12}$.

 



\textbf{Solution 5 (complicated way)}

We can subtract the total areas of triangles DCM and ABC from the rectangle ABCD. For triangle DCM, the base is 4 and the height is 6, so we multiply 4 and 6, then divide by 2 to get 12. For triangle ABC, the base is 4 and the height is 8, so we multiply 4 and 8, then divide by 2 to get 24. We add 24 and 12 to get 36. Then, we calculate the area of rectangle ABCD, which is 48. We subtract 36 from 48, resulting in $\boxed{\textbf{(A) } 12}$.

 -BananaBall00

 


	\item (\textit{AMC 8 2016}) Four students take an exam. Three of their scores are $70$, $80$, and $90$. If the average of their four scores is $70$, then what is the remaining score?


 $\textbf{(A) }40\qquad\textbf{(B) }50\qquad\textbf{(C) }55\qquad\textbf{(D) }60\qquad \textbf{(E) }70$


 
 \textbf{Answer: }A



\textbf{Solution 1}

Let $r$ be the remaining student's score.  We know that the average, 70, is equal to $\frac{70 + 80 + 90 + r}{4}$.  We can use basic algebra to solve for $r$: 
 \[\frac{70 + 80 + 90 + r}{4} = 70\] 
 \[\frac{240 + r}{4} = 70\] 
 \[240 + r = 280\] 
 \[r = 40\] giving us the answer of $\boxed{\textbf{(A)}\ 40}$.

 



\textbf{Solution 2}

Since $90$ is $20$ more than $70$, and $80$ is $10$ more than $70$, for $70$ to be the average, the other number must be $30$ less than $70$, or  $\boxed{\textbf{(A)}\ 40}$.

 


	\item (\textit{AMC 8 2016}) When Cheenu was a boy, he could run $15$ miles in $3$ hours and $30$ minutes. As an old man, he can now walk $10$ miles in $4$ hours. How many minutes longer does it take for him to travel a mile now compared to when he was a boy?


 $\textbf{(A) }6\qquad\textbf{(B) }10\qquad\textbf{(C) }15\qquad\textbf{(D) }18\qquad \textbf{(E) }30$


 
 \textbf{Answer: }B



\textbf{Solution 1}

When Cheenu was a boy, he could run $15$ miles in $3$ hours and $30$ minutes $= 3\times60 + 30$ minutes $= 210$ minutes, thus running $\frac{210}{15} = 14$ minutes per mile. Now that he is an old man, he can walk $10$ miles in $4$ hours $= 4 \times 60$ minutes $= 240$ minutes, thus walking $\frac{240}{10} = 24$ minutes per mile. Therefore, it takes him $\boxed{\textbf{(B)}\ 10}$ minutes longer to walk a mile now compared to when he was a boy.

 



\textbf{Solution 2}

From the question, the old man can travel five miles in two hours, so we can set both speeds to $15$ miles. We can see that Cheenu as an old man takes $2$ hours and $30$ minutes for him to travel $15$ miles, which is also $150$ minutes. We can then divide this by fifteen, which gives us $10$, thus the answer is $\boxed{\textbf{B) 10}}$

 


	\item (\textit{AMC 8 2016}) The number $N$ is a two-digit number.


 • When $N$ is divided by $9$, the remainder is $1$.


 • When $N$ is divided by $10$, the remainder is $3$.


 What is the remainder when $N$ is divided by $11$?


 

$\textbf{(A) }0\qquad\textbf{(B) }2\qquad\textbf{(C) }4\qquad\textbf{(D) }5\qquad \textbf{(E) }7$


 
 \textbf{Answer: }E



\textbf{Solution 1}

From the second bullet point, we know that the second digit must be $3$, for a number divisible by $10$ ends in zero. Since there is a remainder of $1$ when $N$ is divided by $9$, the multiple of $9$ must end in a $2$ for it to have the desired remainder$\pmod {10}.$ We now look for this one: 

 $9(1)=9\\ 9(2)=18\\ 9(3)=27\\ 9(4)=36\\ 9(5)=45\\ 9(6)=54\\ 9(7)=63\\ 9(8)=72$

 The number $72+1=73$ satisfies both conditions. We subtract the biggest multiple of $11$ less than $73$ to get the remainder. Thus, $73-11(6)=73-66=\boxed{\textbf{(E) }7}$.

 



\textbf{Solution 2}

We know that the number has to be one more than a multiple of $9$, because of the remainder of one, and the number has to be $3$ more than a multiple of $10$, which means that it has to end in a $3$. Now, if we just list the first few multiples of $9$ adding one to the number we get: $10, 19, 28, 37, 46, 55, 64, 73, 82, 91$. As we can see from these numbers,  the only one that has a three in the units place is $73$, thus we divide $73$ by $11$, getting $6$ $R7$, hence, $\boxed{\textbf{(E) }7}$.-fn106068

 We could also remember that, for a two-digit number to be divisible by $9$, the sum of its digits has to be equal to $9$. Since the number is one more than a multiple of $9$, the multiple we are looking for has a ones digit of $2$, and therefore a tens digit of $9-2 = 7$, and then we could proceed as above. -vaisri

 


	\item (\textit{AMC 8 2016}) The following bar graph represents the length (in letters) of the names of 19 people. What is the median length of these names?

                    \begin{figure}[H]
                        \centering
                        \includegraphics[width=0.5\textwidth]{Problem/AMC_8_AMC_2016_6_Problem_Figure_1.png}
                    \end{figure}
                    


 $\textbf{(A) }3\qquad\textbf{(B) }4\qquad\textbf{(C) }5\qquad\textbf{(D) }6\qquad \textbf{(E) }7$


 
 \textbf{Answer: }B



\textbf{Solution 1}

We first notice that the median name willl be the $(19+1)/2=10^{\mbox{th}}$ name. The $10^{\mbox{th}}$ name is $\boxed{\textbf{(B)}\ 4}$.

 



\textbf{Solution 2}

To find the median length of a name from a bar graph, we must add up the number of names. Doing so gives us $7 + 3 + 1 + 4 + 4 = 19$. Thus the index of the median length would be the 10th name. Since there are $7$ names with length $3$, and $3$ names with length $4$, the $10$th name would have $4$ letters. Thus our answer is $\boxed{\textbf{(B)}\ 4}$.

 


	\item (\textit{AMC 8 2016}) Which of the following numbers is not a perfect square?


 $\textbf{(A) }1^{2016}\qquad\textbf{(B) }2^{2017}\qquad\textbf{(C) }3^{2018}\qquad\textbf{(D) }4^{2019}\qquad \textbf{(E) }5^{2020}$


 
 \textbf{Answer: }B



\textbf{Solution 1}

Our answer must have an odd exponent in order for it to not be a square.  Because $4$ is a perfect square, $4^{2019}$ is also a perfect square, so our answer is $\boxed{\textbf{(B) }2^{2017}}$.

 



\textbf{Solution 2 (more detailed explanation of Solution 1)}

Any even integer $m$ can be divided into $2n$ where $n$ is an integer as well. Therefore, $\textbf{(A), (C),}$ and $\textbf{(E)}$ are immediately disqualified, as when $m$ is even, $a^m$ can always be written as $(a^\frac{m}{2})^2$ as $(a^m)^n = a^{mn}$ and vice versa. Only odd exponents are non-squares, therefore. $4=2^2$, which means $4^{2019}=(2^2)^2019=(2^{2019})^2$ by the $(a^m)^n = a^{mn}$ rule. The only option left is $\boxed{\textbf{(B) }2^{2017}}$.

 


	\item (\textit{AMC 8 2016}) Find the value of the expression:

 \[100-98+96-94+92-90+\cdots+8-6+4-2.\]$\textbf{(A) }20\qquad\textbf{(B) }40\qquad\textbf{(C) }50\qquad\textbf{(D) }80\qquad \textbf{(E) }100$


 
 \textbf{Answer: }C



\textbf{Solution 1}

We can group each subtracting pair together:
 \[(100-98)+(96-94)+(92-90)+ \ldots +(8-6)+(4-2).\]After subtracting, we have:
 \[2+2+2+\ldots+2+2=2(1+1+1+\ldots+1+1).\]There are $50$ even numbers, therefore there are $\dfrac{50}{2}=25$ even pairs. Therefore the sum is $2 \cdot 25=\boxed{\textbf{(C) }50}$

 



\textbf{Solution 2}

Since our list does not end with one, we divide every number by 2 and we end up with
 \[50-49+48-47+ \ldots +4-3+2-1\]We can group each subtracting pair together:
 \[(50-49)+(48-47)+(46-45)+ \ldots +(4-3)+(2-1).\]There are now $25$ pairs of numbers, and the value of each pair is $1$.  This sum is $25$. However, we divided by $2$ originally so we will multiply $2*25$ to get the final answer of $\boxed{\textbf{(C) }50}$

 



\textbf{Solution 3}

Note: This solution is very similar to Solution 1, but simplified.

 The value of the expression increases by 2 after every pair as shown:

 Pair 1: [100-98= $2$]

 Pair 1 + 2: [2+96-94= $4$]

 Pair 1 + 2 + 3: [4+92-90= $6$]

 Since there is 25 pairs in the expression, the value will increase by two 25 times. Thus, the sum of the expression is  $2 \cdot 25=\boxed{\textbf{(C) }50}$

 


	\item (\textit{AMC 8 2016}) What is the sum of the distinct prime integer divisors of $2016$?


 $\textbf{(A) }9\qquad\textbf{(B) }12\qquad\textbf{(C) }16\qquad\textbf{(D) }49\qquad \textbf{(E) }63$


 
 \textbf{Answer: }B



\textbf{Solution 1}

The prime factorization is $2016=2^5\times3^2\times7$.  Since the problem is only asking us for the distinct prime factors, we have $2,3,7$.  Their desired sum is then $\boxed{\textbf{(B) }12}$.

 


	\item (\textit{AMC 8 2016}) Suppose that $a * b$ means $3a-b.$ What is the value of $x$ if

 \[2 * (5 * x)=1\]
$\textbf{(A) }\frac{1}{10} \qquad\textbf{(B) }2\qquad\textbf{(C) }\frac{10}{3} \qquad\textbf{(D) }10\qquad \textbf{(E) }14.$


 
 \textbf{Answer: }D



\textbf{Solution 1}

Let us plug in $(5 * x)=1$ into $3a-b$. Thus it would be $3(5)-x$. Now we have $2*(15-x)=1$. Plugging $2*(15-x)$ into $3a-b$, we have $6-15+x=1$. Solving for $x$ we have 
 \[-9+x=1\]
 \[x=\boxed{\textbf{(D)}   \,        10}\]

 



\textbf{Solution 2}

Let us set a variable $y$ equal to $5 * x$. Solving for y in the equation $3(2)-y=1$, we see that y is equal to five. By substitution, we see that $5 * x$ = 5. Solving for x in the equation $5(3)-x = 5$ we get 
 \[x=\boxed{\textbf{(D)}   \,        10}\]

 


	\item (\textit{AMC 8 2016}) Determine how many two-digit numbers satisfy the following property: when the number is added to the number obtained by reversing its digits, the sum is $132.$


 $\textbf{(A) }5\qquad\textbf{(B) }7\qquad\textbf{(C) }9\qquad\textbf{(D) }11\qquad \textbf{(E) }12$


 
 \textbf{Answer: }B



\textbf{Solution 1}

We can see that the original number can be written as $10a+b$, where $a$ represents the tens digit and $b$ represents the units digit. When this number is added to the number obtained by reversing its digits, which is $10b+a$, the sum would be $11a+11b$. From this, we can construct the equation $11a+11b=132$, which simplifies to $a+b=12$. Since there are 7 pairs of such digits $a$ and $b$, $(3,9), (4,8), (5,7), (6,6), (7,5), (8,4), (9,3)$, the answer would be $\boxed{\textbf{(B) } 7}.$

 



\textbf{Solution 2}

We can set the number as ab where a is the tens digit and b is the ones digit. So now the equation will be ab+ba=132. a+b has to have a remainder of 2 when divided by 10 so it will be \(a+b\equiv 2\quad (\bmod 10)\). We also know a<10 and b<10. So a+b can either be 2 or 12. a+b cannot be 2 because then there will be only 3 numbers that work and that isn't in the answer choice. So a+b=12 . To check this we can do ab+ba=132 which equals to (a+b)0+(a+b)=132 and since we said a+b=12 we get 120+12=132 which is true. So we have a+b=12 and a<10 and b<10.If a is 9 then b=3 and if a=3 then b=9 calculating how many pairs are in between you get 7. So the answer is $\boxed{\textbf{(B) } 7}.$

 



\textbf{Solution 3}

Like similar solutions, we can say that \( 10a+b \) is a two digit number, and therefore we have \( 10a+b + 10b+a = 132 \) which implies \( 11a + 11b = 132 \) which implies \( a + b = 12 \).

 We then use stars and bars to find out that there are \( \binom{13}{1} = 13 \) different possible sums.

 However, we want to satisfy \( 0 \le a \le 9 \) and \( 0 \le b \le 9 \), therefore we see that the sums \( 0+12 \), \( 1+11 \), \( 2+10 \), and \( 12+0 \), \( 11+1 \), and \( 10+2 \) do not work, and our answer is \( 13-6= \) $\boxed{\textbf{(B) } 7}.$.

 



\textbf{Solution 4}

The two digit number 10A+B + the reversed number 10B+A is equal to 11A+11B = 132. Dividing both sides by 11 gives A + B = 12. Since we can't have $1 + 11, 2 + 10$ etc. meaning there are $\boxed{\textbf{(B) } 7}.$ solutions.

 


	\item (\textit{AMC 8 2016}) Jefferson Middle School has the same number of boys and girls. $\frac{3}{4}$ of the girls and $\frac{2}{3}$
of the boys went on a field trip. What fraction of the students on the field trip were girls?


 $\textbf{(A) }\frac{1}{2}\qquad\textbf{(B) }\frac{9}{17}\qquad\textbf{(C) }\frac{7}{13}\qquad\textbf{(D) }\frac{2}{3}\qquad \textbf{(E) }\frac{14}{15}$


 
 \textbf{Answer: }B



\textbf{Solution 1}

Let there be $b$ boys and $g$ girls in the school. We see $g=b$, which means $\frac{3}{4}b+\frac{2}{3}b=\frac{17}{12}b$ kids went on the trip and $\frac{3}{4}b$ kids are girls. So, the answer is $\frac{\frac{3}{4}b}{\frac{17}{12}b}=\frac{9}{17}$, which is $\boxed{\textbf{(B)} \frac{9}{17}}$.

 



\textbf{Solution 2}

Using WLOG (Without loss of generativity), Let there be $12$ boys and $12$ girls in the school. Now we can do $\frac{3}{4}\times{12}$ + $\frac{2}{3}\times{12}$ to get the total number of students going to the field trip to be $17$. Since we already know the number of girls to be $9$. We have our answer to be $\frac{9}{17}$. So, the answer is $\boxed{\textbf{(B)} \frac{9}{17}}$.

 


	\item (\textit{AMC 8 2016}) Two different numbers are randomly selected from the set $\{ - 2, -1, 0, 3, 4, 5\}$ and multiplied together. What is the probability that the product is $0$?


 $\textbf{(A) }\dfrac{1}{6}\qquad\textbf{(B) }\dfrac{1}{5}\qquad\textbf{(C) }\dfrac{1}{4}\qquad\textbf{(D) }\dfrac{1}{3}\qquad \textbf{(E) }\dfrac{1}{2}$


 
 \textbf{Answer: }D



\textbf{Solution 1}

1. Identify the total number of ways to select two different numbers from the set:

 The set has 6 elements. The number of ways to choose 2 different numbers from 6 is given by the combination formula: $\binom{6}{2} = \frac{6 \times 5}{2 \times 1} = 15$.

 2. Identify the favorable outcomes:

 For the product to be zero, one of the chosen numbers must be zero. The set contains one zero (0). To have a product of zero, we need to choose 0 and any other number from the remaining five numbers $-2, -1, 3, 4, 5$.

 The number of ways to choose 0 and one other number from the remaining five is 5.

 3. Calculate the probability:

 The probability is the number of favorable outcomes divided by the total number of outcomes: $\text{Probability} = \frac{\text{Number of favorable outcomes}}{\text{Total number of outcomes}} = \frac{5}{15} = \frac{1}{3}$.

 Thus, the probability that the product is $0$ is $\boxed{\textbf{(D)} \ \frac{1}{3}}.$

 



\textbf{Solution 2 (Complementary Counting)}

Because the only way the product of the two numbers is $0$ is if one of the numbers we choose is $0,$ we calculate the probability of NOT choosing a $0.$ We get $\frac{5}{6} \cdot \frac{4}{5} = \frac{2}{3}.$ Therefore our answer is $1 - \frac{2}{3} = \boxed{\textbf{(D)} \ \frac{1}{3}}.$

 



\textbf{Solution 3(One Sentence!!!)}

1/6 * Any other number + Any other number * 1/6 = 1/6+1/6 = $\boxed{\textbf{(D)} \ \frac{1}{3}}$

 -JasonDaGoat

 
1/6 Chance of 0 * A number in the set + 1/6 chance of A number in the set * 0 = 1/6+1/6 = $\boxed{\textbf{(D)} \ \frac{1}{3}}$.

 - edited by TTDBJ

 



\textbf{Solution 4}

To get a product of $0$, one of the numbers must be $0$, which means that we have $2P1=2$ possible combinations of a random number from the set, and $0$. There is $6$ total possible products, giving us $\boxed{\frac{1}{3}}$.-Electric\_Cheetah\_88

 


	\item (\textit{AMC 8 2016}) Karl's car uses a gallon of gas every $35$ miles, and his gas tank holds $14$ gallons when it is full. One day, Karl started with a full tank of gas, 
drove $350$ miles, bought $8$ gallons of gas, and continued driving to his destination. When he arrived, his gas tank was half full. How many miles did Karl drive that day? 


 $\textbf{(A)}\mbox{ }525\qquad\textbf{(B)}\mbox{ }560\qquad\textbf{(C)}\mbox{ }595\qquad\textbf{(D)}\mbox{ }665\qquad\textbf{(E)}\mbox{ }735$


 
 \textbf{Answer: }A



\textbf{Solution}

Since he uses a gallon of gas every $35$ miles, he had used $\frac{350}{35} = 10$ gallons after $350$ miles. Therefore, after the first leg of his trip he had $14 - 10 = 4$ gallons of gas left. Then, he bought $8$ gallons of gas, which brought him up to $12$ gallons of gas in his gas tank. When he arrived, he had $\frac{1}{2} \cdot 14 = 7$ gallons of gas. So he used $5$ gallons of gas on the second leg of his trip. Therefore, the second part of his trip covered $5 \cdot 35 = 175$ miles. Adding this to the $350$ miles, we see that he drove $350 + 175 = \boxed{\textbf{(A)}    \, 525}$ miles.

 


	\item (\textit{AMC 8 2016}) What is the largest power of $2$ that is a divisor of $13^4 - 11^4$?


 $\textbf{(A)}\mbox{ }8\qquad \textbf{(B)}\mbox{ }16\qquad \textbf{(C)}\mbox{ }32\qquad \textbf{(D)}\mbox{ }64\qquad \textbf{(E)}\mbox{ }128$


 
 \textbf{Answer: }C



\textbf{Solution 1}

First, we use difference of squares on $13^4 - 11^4 = (13^2)^2 - (11^2)^2$ to get $13^4 - 11^4 = (13^2 + 11^2)(13^2 - 11^2)$. Using difference of squares again and simplifying, we get $(169 + 121)(13+11)(13-11) = 290 \cdot 24 \cdot 2 = (2\cdot 8 \cdot 2) \cdot (3 \cdot 145)$. Realizing that we don't need the right-hand side because it doesn't contain any factor of 2, we see that the greatest power of $2$ that is a divisor $13^4 - 11^4$ is $\boxed{\textbf{(C)}\ 32}$.

 



\textbf{Solution 2 (variant of Solution 1)}

Just like in the above solution, we use the difference-of-squares factorization, but only once to get $13^4-11^4=(13^2-11^2)(13^2+11^2).$ We can then compute that this is equal to $48\cdot290.$ Note that $290=2\cdot145$ (we don't need to factorize any further as $145$ is already odd) thus the largest power of $2$ that divides $290$ is only $2^1=2,$ while $48=2^4\cdot3,$ so the largest power of $2$ that divides $48$ is $2^4=16.$ Hence, the largest power of $2$ that is a divisor of $13^4-11^4$ is $2\cdot16=\boxed{\textbf{(C)}~32}.$

 



\textbf{Solution 3 (Lifting the exponent)}

Let $n=13^4-11^4.$ We wish to find the largest power of $2$ that divides $n$.

 Denote $v_p(k)$ as the largest exponent of $p$ in the prime factorization of $n$. In this problem, we have $p=2$. 

 By the Lifting the Exponent Lemma on $n$, 

 
 \[v_2(13^4-11^4)=v_2(13-11)+v_2(4)+v_2(13+11)-1\]
 \[=v_2(2)+v_2(4)+v_2(24)-1\]
 \[=1+2+3-1=5.\]

 Therefore, exponent of the largest power of $2$ that divids $13^4-11^4$ is $5,$ so the largest power of $2$ that divides this number is $2^5=\boxed{\textbf{(C)} 32}$.

 



\textbf{Solution 4 (Brute Force)}

We can simply take 13 to the 4th power, which is 28561. We subtract that by 11 to the 4th power, which is 14641 (You can use Pascal's Triangle to find this). Finally, subtract the numbers to get 13920. 

 To test the options, since we need the largest one, we can go from top down. Testing, we see that both D and E are decimals,and 32 works. So, our answer is $\boxed{\textbf{(C)}~32}.$

 -themathgood

 
Note: This is not the fastest way, and is not recommended.

 



\textbf{Solution 5}

$13^4 - 11^4$ is $(13^2-11^2)(13^2+11^2)$ which is $(169-121)(169+121)$ which is $(48)(290)$ which has 5 factors of 2 making $\boxed{\textbf{(C)}~32}.$

 


	\item (\textit{AMC 8 2016}) Annie and Bonnie are running laps around a $400$-meter oval track. They started at the same time, at the same place, but Annie has pulled ahead because she runs $25\%$ faster than Bonnie. How many laps will Annie have run when she first passes Bonnie?


 $\textbf{(A) }1\dfrac{1}{4}\qquad\textbf{(B) }3\dfrac{1}{3}\qquad\textbf{(C) }4\qquad\textbf{(D) }5\qquad \textbf{(E) }25$


 
 \textbf{Answer: }D



\textbf{Solution 1}

Each lap Bonnie runs, Annie runs another quarter lap, so Bonnie will run four laps before she is overtaken because it takes 4 of $\frac{1}{4}$ of the total lap. to total one lap. When Bonnie runs 4 laps, Annie has gained one full lap and would overtake Bonnie. That means Annie will have run $\boxed{\textbf{(D)}\ 5 }$ laps.

 



\textbf{Solution 2}

Saying that Bonnie runs 40 meters per second, then Annie will run 50 meters per second. It will take Bonnie 10 seconds to run 1 lap and it would take 8 seconds for Annie to do 1 lap. The LCM of 8 and 10 is 40. In 40 seconds Bonnie will do 4 laps and Annie will do 5 laps. Since Annie did 1 more lap she passed Bonnie. So the answer is $\boxed{\textbf{(D)}\ 5 }$ laps.

 


	\item (\textit{AMC 8 2016}) An ATM password at Fred's Bank is composed of four digits from $0$ to $9$, with repeated digits allowable. If no password may begin with the sequence $9,1,1,$ then how many passwords are possible?


 $\textbf{(A)}\mbox{ }30\qquad\textbf{(B)}\mbox{ }7290\qquad\textbf{(C)}\mbox{ }9000\qquad\textbf{(D)}\mbox{ }9990\qquad\textbf{(E)}\mbox{ }9999$


 
 \textbf{Answer: }D



\textbf{Solution 1}

For the first three digits, there are $10^3-1=999$ combinations since $911$ is not allowed. For the final digit, any of the $10$ numbers are allowed. $999 \cdot 10 = 9990 \rightarrow \boxed{\textbf{(D)}\ 9990}$.

 



\textbf{Solution 2}

Counting the prohibited cases, we find that there are 10 of them. This is because, when we start with 9,1, and 1, we can have any of the 10 digits for the last digit. So, our answer is $10^4-10=\boxed{\textbf{(D)}\ 9990}.$

 


	\item (\textit{AMC 8 2016}) In an All-Area track meet, $216$ sprinters enter a $100-$meter dash competition. The track has $6$ lanes, so only $6$ sprinters can compete at a time. At the end of each race, the five non-winners are eliminated, and the winner will compete again in a later race. How many races are needed to determine the champion sprinter?


 $\textbf{(A)}\mbox{ }36\qquad\textbf{(B)}\mbox{ }42\qquad\textbf{(C)}\mbox{ }43\qquad\textbf{(D)}\mbox{ }60\qquad\textbf{(E)}\mbox{ }72$


 
 \textbf{Answer: }C



\textbf{Solution 1}

From any $n-$th race, only $\frac{1}{6}$ will continue on. Since we wish to find the total number of races, a column representing the races over time is ideal.Starting with the first race:
 \[\frac{216}{6}=36\]
 \[\frac{36}{6}=6\]
 \[\frac{6}{6}=1\]Adding all of the numbers in the second column yields $\boxed{\textbf{(C)}\ 43}$

 



\textbf{Solution 2}

Every race eliminates $5$ players. The winner is decided when there is only $1$ runner left. You can construct the equation: $216$ - $5x$ = $1$. Thus, $215$ players have to be eliminated. Therefore, we need $\frac{215}{5}$ games to decide the winner, or $\boxed{\textbf{(C)}\ 43}$

 


	\item (\textit{AMC 8 2016}) The sum of $25$ consecutive even integers is $10,000$. What is the largest of these $25$ consecutive integers?


 $\textbf{(A)}\mbox{ }360\qquad\textbf{(B)}\mbox{ }388\qquad\textbf{(C)}\mbox{ }412\qquad\textbf{(D)}\mbox{ }416\qquad\textbf{(E)}\mbox{ }424$


 
 \textbf{Answer: }E



\textbf{Solution 1}

Let $n$ be the 13th consecutive even integer that's being added up. By doing this, we can see that the sum of all 25 even numbers will simplify to $25n$ since $(n-2k)+\dots+(n-4)+(n-2)+(n)+(n+2)+(n+4)+ \dots +(n+2k)=25n$. Now, $25n=10000 \rightarrow n=400$. Remembering that this is the 13th integer, we wish to find the 25th, which is $400+2(25-13)=\boxed{\textbf{(E)}\ 424}$.

 



\textbf{Solution 2}

Let $x$ be the smallest number. The equation will become, $x+(x+2)+(x+4)+\cdots +(x+48)=10,000$. After you combine like terms, you get $25x+(50*12)=10,000$ which turns into $10,000-600=25x$. $25x=9400$, so $x=376$. Then, you add $376+48 = \boxed{\textbf{(E)}\ 424}$.

 Edit(dnathan12): For clarity, you add 48, because you want the largest number in the list. 376 is the smallest number~AfterglowBlaziken

 


	\item (\textit{AMC 8 2016}) The least common multiple of $a$ and $b$ is $12$, and the least common multiple of $b$ and $c$ is $15$. What is the least possible value of the least common multiple of $a$ and $c$?


 $\textbf{(A) }20\qquad\textbf{(B) }30\qquad\textbf{(C) }60\qquad\textbf{(D) }120\qquad \textbf{(E) }180$


 
 \textbf{Answer: }A



\textbf{Solution}

We wish to find possible values of $a$, $b$, and $c$. By finding the greatest common factor of $12$ and $15$, we can find that $b$ is 3. Moving on to $a$ and $c$, in order to minimize them, we wish to find the least such that the least common multiple of $a$ and $3$ is $12$, $\rightarrow 4$. Similarly, with $3$ and $c$, we obtain $5$. The least common multiple of $4$ and $5$ is $20 \rightarrow \boxed{\textbf{(A)} 20}$

 



\textbf{Solution 2}

The factors of 12 are 12,6,4,3,2,1. The factors of 15 are 1,3,5,15. The 2 numbers that repeat are 1 and 3 so b either has to be 1 or 3. If b is 3 then a is 4 and c is 5 and the least common multiple of 4 and 5 are 20. We don't have to test 1 because 20 is the lowest answer, so if b equaling 1 resulted in the least common multiple being less that 20 then the correct answer won't be there. So the answer is $\boxed{\textbf{(A)} 20}$.

 



\textbf{Solution 3 (Trial & Error)}

Let's say that $a = 4, b = 3, c = 5$ (because $LCM(3,4)=12$ and $LCM(3,5)=15$), this works so clearly $GCF(4, 5) = \boxed{\textbf{(A)} 20}$. ~ JasonDaMoat, ~ Minor edited by Quetzalcoatlus

 


	\item (\textit{AMC 8 2016}) A top hat contains 3 red chips and 2 green chips. Chips are drawn randomly, one at a time without replacement, until all 3 of the reds are drawn or until both green chips are drawn. What is the probability that the 3 reds are drawn?


 $\textbf{(A) }\dfrac{3}{10}\qquad\textbf{(B) }\dfrac{2}{5}\qquad\textbf{(C) }\dfrac{1}{2}\qquad\textbf{(D) }\dfrac{3}{5}\qquad \textbf{(E) }\dfrac{7}{10}$


 
 \textbf{Answer: }B



\textbf{Solution 1}

We put five chips randomly in order and then pick the chips from the left to the right. To find the number of ways to rearrange the three red chips and two green chips, we solve for $\binom{5}{2} = 10$. However, we notice that whenever the last chip we draw is red, we pick both green chips before we pick the last red chip. Similarly, when the last chip is green, we pick all three red chips before the last green chip. This means that the last chip must be green in all the situations that work. This means we are left with finding the number of ways to rearrange three red chips and one green chip, which is $\binom{4}{3} = 4$. Because a green chip will be last $4$ out of the $10$ situations, our answer is $\boxed{\textbf{(B) } \frac{2}{5}}$.

 



\textbf{Solution 2}

We can use basic probability to find the answer to this question. The probability of the first chip being red is $\frac{3}{5}$ because there are $3$ red chips and $5$ chips in total. The probability of the second chip being red is $\frac{2}{4}$ because there are now $2$ red chips and $4$ total chips. The probability of the third chip being red is $\frac{1}{3}$ because there are now $1$ red chips and $3$ total chips. If we multiply these fractions, we get $\frac{1}{10}$. To get the answer of $\frac{2}{5}$ we need to multiply this product by $4$ because there are $4$ different draws that these red chips can be chosen on. They can be chosen on the $\text{1st 2nd, 3rd, or 4th}$ draws because if a red chip was chosen on the $\text{5th}$ draw, $2$ green chips would have already been chosen. So our answer is $\frac{1}{10}$ multiplied by $4$ which would leave us with the answer of $\boxed{\textbf{(B) } \frac{2}{5}}$.

 \textbf{Faster Solution (Similar to Solution 1)} Randomly arrange them in a row of 5. There is a 3/5 chance that the last one will be red. If the last one is red, 3 reds will not be drawn first, so there is a $\boxed{\textbf{(B) } \frac{2}{5}}$ chance that 3 reds will come out first.

 


	\item (\textit{AMC 8 2016}) Rectangle $DEFA$ below is a $3 \times 4$ rectangle with $DC=CB=BA=1$. What is the area of the "bat wings" (shaded region)?
                    \begin{figure}[H]
                        \centering
                        \includegraphics[width=0.5\textwidth]{Problem/AMC_8_AMC_2016_22_Problem_Figure_1.png}
                    \end{figure}
                    


 $\textbf{(A) }2\qquad\textbf{(B) }2 \frac{1}{2}\qquad\textbf{(C) }3\qquad\textbf{(D) }3 \frac{1}{2}\qquad \textbf{(E) }4$


 
 \textbf{Answer: }C



\textbf{Solution 1}

The area of trapezoid $CBFE$ is $\frac{1+3}2\cdot 4=8$. Next, we find the height of each triangle to calculate their area. The two non-colored isosceles triangles are similar, and are in a $3:1$ ratio by AA similarity (alternate interior and vertical angles) so the height of the larger is $3,$ while the height of the smaller one is $1.$ Thus, their areas are $\frac12$ and $\frac92$. Subtracting these areas from the trapezoid, we get $8-\frac12-\frac92 =\boxed3$. Therefore, the answer to this problem is $\boxed{\textbf{(C) }3}$.

 



\textbf{Solution 2}

Plot the point $G$ where the two "wings" intersect. Now notice how $\triangle CBG\sim\triangle EFG$. Since the length of $\overline {CB}$ is one third that of $\overline {EF}$, then that means $\triangle EFG$'s height is $3$ times bigger than  $\triangle CBG$. Since both of their heights ($h$) add up to four, then we have the equation $3h+h=4 \implies h=1$. Now that we now the height and length of both triangles, we can use complementary counting, $\text{Area}-\text{Unshaded Region}$.

 $\text {Total Area}=12$

 $[\triangle CDE]=2$

 $[\triangle ABF]=2$

 $[\triangle CBG]=\frac1{2}$

 $[\triangle EFG]=\frac{9}{2}$

 $\text {Unshaded Region}=9\implies\text{"Bat Wings"}=\boxed{\textbf{(C) }3}$

 


	\item (\textit{AMC 8 2016}) Two congruent circles centered at points $A$ and $B$ each pass through the other circle's center. The line containing both $A$ and $B$ is extended to intersect the circles at points $C$ and $D$. The circles intersect at two points, one of which is $E$. What is the degree measure of $\angle CED$?


 $\textbf{(A) }90\qquad\textbf{(B) }105\qquad\textbf{(C) }120\qquad\textbf{(D) }135\qquad \textbf{(E) }150$


 
 \textbf{Answer: }C



\textbf{Solution 1}

Observe that $\triangle{EAB}$ is equilateral (all are radii of congruent circles). Therefore, $m\angle{AEB}=m\angle{EAB}=m\angle{EBA} = 60^{\circ}$. Since $CD$ is a straight line, we conclude that $m\angle{EBD} = 180^{\circ}-60^{\circ}=120^{\circ}$. Since $BE=BD$ (both are radii of the same circle), $\triangle{BED}$ is isosceles, meaning that $m\angle{BED}=m\angle{BDE}=30^{\circ}$. Similarly, $m\angle{AEC}=m\angle{ACE}=30^{\circ}$. 

 Now, $\angle{CED}=m\angle{AEC}+m\angle{AEB}+m\angle{BED} = 30^{\circ}+60^{\circ}+30^{\circ} = 120^{\circ}$. Therefore, the answer is $\boxed{\textbf{(C) }\ 120}$.

 \textbf{Graph for Solution 1} \href{https://drive.google.com/file/d/1sNkKxIi7\_SpxihTp0QrOcD3Js14buvd8/view?usp=drivesdk}{https://drive.google.com/file/d/1sNkKxIi7\_SpxihTp0QrOcD3Js14buvd8/view?usp=drivesdk}~ By Quetzalcoatlus

 


	\item (\textit{AMC 8 2016}) The digits $1$, $2$, $3$, $4$, and $5$ are each used once to write a five-digit number $PQRST$. The three-digit number $PQR$ is divisible by $4$, the three-digit number $QRS$ is divisible by $5$, and the three-digit number $RST$ is divisible by $3$. What is $P$?


 $\textbf{(A) }1\qquad\textbf{(B) }2\qquad\textbf{(C) }3\qquad\textbf{(D) }4\qquad \textbf{(E) }5$


 
 \textbf{Answer: }A



\textbf{Solution 1 (Modular Arithmetic)}

We see that since $QRS$ is divisible by $5$, $S$ must equal either $0$ or $5$, but it cannot equal $0$, so $S=5$. We notice that since $PQR$ must be even, $R$ must be either $2$ or $4$. However, when $R=2$, we see that $T \equiv 2 \pmod{3}$, which cannot happen because $2$ and $5$ are already used up; so $R=4$. This gives $T \equiv 3 \pmod{4}$, meaning $T=3$. Now, we see that $Q$ could be either $1$ or $2$, but $14$ is not divisible by $4$, but $24$ is. This means that $Q=2$ and $P=\boxed{\textbf{(A)}\ 1}$.

 Editor's note: If you are in a time constraint and can't get through the whole problem, you can, with this modular arithmetic,  eliminate the ones' you know can't be in the first position.

 



\textbf{Solution 2}

We know that out of $PQRST,$ $QRS$ is divisible by $5$. Therefore $S$ is obviously 5 because $QRS$ is divisible by 5. So we now have $PQR5T$ as our number. Next, let's move on to the second piece of information that was given to us. $RST$ is divisible by 3. So, according to the divisibility by 3 rule, the sum of $RST$ has to be a multiple of 3. The only 2 big enough are 9 and 12 and since 5 is already given. The possible sums of $RT$ are 4 and 7. So, the possible values for $R$ are 1,3,4,3 and the possible values of $T$ are 3,1,3,4. So, using this we can move on to the fact that $PQR$ is divisible by 4. So, using that we know that $R$ has to be even so 4 is the only possible value for $R$. Using that we also know that 3 is the only possible value for 3. So, we have $PQRST$ = $PQ453$ so the possible values are 1 and 2 for $P$ and $Q$. Using the divisibility rule of 4 we know that $QR$ has to be divisible by 4. So, either 14 or 24 are the possibilities, and 24 is divisible by 4. So the only value left for $P$ is 1.  $P=\boxed{\textbf{(A)}\ 1}$.

 



\textbf{Solution 3}

We can simply try each of the answer choices, and we will see which one works. Trying $P=\boxed{\textbf{(A) }1}$, if $PQR$ is divisible by $4$, $QR$ must be divisible by four. Therefore, $QR$ can only be $24$, $52$, or $32$. However, since $QRS$ is divisible by $5$, $S = 5$, so $QR$ cannot be $52$. When $QR = 32$, $R = 2$, the last requirement cannot be satisfied because $R + S + T = 2 + 4 + 5 = 11$, and $11$ is not divisible by $3$. However, when $QR = 24$, $R = 4$, the last requirement can be satisfied. Hence, we can see that when $P=\boxed{\textbf{(A) }1}$, there is one way to satisfy all three requirements, leading to a conclusion that $P$ is $\boxed{\textbf{(A) }1}$.

 


	\item (\textit{AMC 8 2016}) A semicircle is inscribed in an isosceles triangle with base $16$ and height $15$ so that the diameter of the semicircle is contained in the base of the triangle as shown. What is the radius of the semicircle?


 
                    \begin{figure}[H]
                        \centering
                        \includegraphics[width=0.5\textwidth]{Problem/AMC_8_AMC_2016_25_Problem_Figure_1.png}
                    \end{figure}
                    


 $\textbf{(A) }4 \sqrt{3}\qquad\textbf{(B) } \dfrac{120}{17}\qquad\textbf{(C) }10\qquad\textbf{(D) }\dfrac{17\sqrt{2}}{2}\qquad \textbf{(E)} \dfrac{17\sqrt{3}}{2}$


 
 \textbf{Answer: }B



\textbf{Solution 1}


                    \begin{figure}[H]
                        \centering
                        \includegraphics[width=0.5\textwidth]{Solution/AMC_8_AMC_2016_25_Solution_Figure_1.png}
                    \end{figure}
                    

 First, we drop a perpendicular, shown above, to the base of the triangle, cutting the triangle into two congruent right triangles. This triangle is isosceles, which means perpendiculars are medians and vice versa. The base of the resulting right triangle is $8$ for both sides, and the height is $15,$ as given. Using the Pythagorean theorem, we can find the length of the hypotenuse, or $17.$ Using the two legs of the right triangle, we find the area of the right triangle, $60$. $\frac{60}{17}$ times $2$ results in the radius, which is the height of the right triangle when using the hypotenuse as the base. Hence, the answer is $\boxed{\textbf{(B) }\frac{120}{17}}$.

 
                \begin{itemize}
                    \item Note — There are several other following solutions below that use similar methods of finding the area two different ways shown in this solution. Try to eliminate any words (such as 'can' or 'would be') that make your solution less concise. Try to write 'sure of yourself.'
                \end{itemize}


                



\textbf{Solution 2: Similar Triangles}


                    \begin{figure}[H]
                        \centering
                        \includegraphics[width=0.5\textwidth]{Solution/AMC_8_AMC_2016_25_Solution_Figure_2.png}
                    \end{figure}
                    Let's call the triangle $\triangle ABC,$ where $AB=16$ and $AC=BC.$ Let's say that $D$ is the midpoint of $AB$ and $E$ is the point where $AC$ is tangent to the semicircle. We could also use $BC$ instead of $AC$ because of symmetry.

 Notice that $\triangle ACD \cong \triangle BCD,$ and are both 8-15-17 right triangles. We also know that we create a right angle with the intersection of the radius and a tangent line of a circle (or part of a circle). So, by $AA$ similarity, $\triangle AED \sim \triangle ADC,$ with $\angle EAD \cong \angle DAC$ and $\angle CDA \cong \angle DEA.$ This similarity means that we can create a proportion: $\frac{AD}{AC}=\frac{DE}{CD}.$ We plug in $AD=\frac{AB}{2}=8, AC=17,$ and $CD=15.$ After we multiply both sides by $15,$ we get $DE=\frac{8}{17} \cdot 15= \boxed{\textbf{(B) }\frac{120}{17}}.$

 (By the way, we could also use $\triangle DEC \sim \triangle ADC.$)

 



\textbf{Solution 3: Inscribed Circle}


                    \begin{figure}[H]
                        \centering
                        \includegraphics[width=0.5\textwidth]{Solution/AMC_8_AMC_2016_25_Solution_Figure_3.png}
                    \end{figure}
                    

 We'll call this triangle $\triangle ABD$. Let the midpoint of base $BD$ be $M$. Divide the triangle in half by drawing a line from $A$ to $M$. Half the base of $\triangle ABD$ is $\frac{16}{2} = 8$. The height is $15$, which is given in the question. Using the Pythagorean Triple $8$-$15$-$17$, the length of each of the legs ($AB$ and $DA$) is 17.

 Reflect the triangle over its base. This will create an inscribed circle in a rhombus $ABCD$. Because $AB \cong DA$, $BC \cong CD$. Therefore $AB = BC = CD = DA$.

 The semiperimeter $s$ of the rhombus is $\frac{AB + BC + CD + DA}{2} = \frac{(17)(4)}{2} = 34$. Since the area of $\triangle ABD$ is $\frac{bh}{2}$, the area $[ABCD]$ of the rhombus is twice that, which is $bh = (16)(15) = 240$.

 The \href{https://en.wikipedia.org/wiki/Incircle\_and\_excircles\_of\_a\_triangle\#Incircle}{Formula for the Incircle of a Quadrilateral} is $s$$r$ = $[ABCD]$. Substituting the semiperimeter and area into the equation, $34r = 240$. Solving this, $r = \frac{240}{34}$ = $\boxed{\textbf{(B) }\frac{120}{17}}$.

 



\textbf{Solution 4: Inscribed Circle}

Noting that we have a 8-15-17 triangle, we can find $AE$ and $CE.$ Let $AE=x$, $CE=17-x.$ Then by similar triangles (or "Altitude on Hypotenuse") we have $15^2=x*17.$ Thus, $AE=x=225/17, CE=64/17.$ Now again by "Altitude on Hypotenuse”, $r=\sqrt{AE*CE}.$ Therefore $r=\boxed{\textbf{(B) }\frac{120}{17}}$.

 



\textbf{Solution 5: Simple Trigonometry}

Note: This solution uses \href{https://artofproblemsolving.comhttp://artofproblemsolving.com/wiki/index.php?title=Trigonometry\#Word}{Trigonometric Concepts}
                    \begin{figure}[H]
                        \centering
                        \includegraphics[width=0.5\textwidth]{Solution/AMC_8_AMC_2016_25_Solution_Figure_4.png}
                    \end{figure}
                    

 Denote the bottom left vertex of the isosceles triangle to be $8$

 Denote the bottom right vertex of the isosceles triangle to be $8$ 

 Denote the top vertex of the isosceles triangle to be $x$

 Drop an altitude from $x$ to side $AB$. Denote the foot of the intersection to be $D$.

 By the Pythagorean Theorem, $AC=17$.

 Now, we see that by sin(x), $\sin{A}=\frac{15}{17}$.

 This implies that $\sin{A}=\frac{r}{8}$ (r=radius of semicircle).

 Hence, $r=\boxed{\textbf{(B) }\frac{120}{17}}$.

 



\textbf{Solution 6: Area}


                    \begin{figure}[H]
                        \centering
                        \includegraphics[width=0.5\textwidth]{Solution/AMC_8_AMC_2016_25_Solution_Figure_2.png}
                    \end{figure}
                    Credits for Asymptote go to whoever wrote this diagram up in Solution 2. 

 There are two ways to find the area of $\triangle ABC$. The first way is the most obvious, and that is to multiply the base times the height ($16\cdot15$) and then divide it by two. The second way is, in a way, a little more complex. Note that $\triangle ACD$ and $\triangle BCD$ are congruent. This means that if we find the area of one triangle, we can just multiply it's area by 2 and we've found the area of the larger triangle $ABC$. But since we always divide by two, as it is in the formula for finding the area of a triangle, the multiply by two and divide by two cancel out, giving us that if we just multiply base times height of $\triangle ACD$, we will get the area of $\triangle ABC$.

 Now you might think: "But what is the other way of finding the area". Well that is $AC$, which would be the base, times $DE$, the radius, which would be the height. 

 The first way to find the area gives us the area of the $\triangle ABC$, $\frac{15\cdot16}{2}=120$. This gives us $120=AC\cdot r$ ($r$ signifies radius). We can find $AC$ using the Pythagorean Theorem on $\triangle ACD$. The two legs are $8$ and $15$, which gives us that the hypotenuse, $AC$, is equal to $17$. 

 Now that we have an equation with only one variable for the radius, $120=17r$, we can just solve for the radius. We get $r=\boxed{\textbf{(B) }\frac{120}{17}}$. 

 This may seem like a lengthy explanation, but doing it yourself when you know what to do, it actually takes very little time. Try it yourself!

 



\textbf{Solution 7}

Let us draw altitude $\overline{CD}.$ This cuts our base into line segments with length $\dfrac{16}{2}=8.$ Finding the area of the resulting triangles gives $[\triangle ADC] = [\triangle BDC] = \dfrac{8 \cdot 15}{2} = 60.$  Since $m\angle CDB = m\angle CDA = 90^\circ,$ we use the Pythagorean Theorem to find length $\overline{AC}$:
 \begin{align*} (\overline{AD})^2+(\overline{CD})^2 &= (\overline{AC})^2 \\ 8^2+15^2 &= (\overline{AC})^2 \\ 64+225 &= (\overline{AC})^2 \\ 289 &= (\overline{AC})^2 \\ \sqrt{289} &= \sqrt{(\overline{AC})^2} \\ 17 &= AC. \end{align*}
Thus we have $\overline{AC}=\overline{BC}=17.$ 

 
                    \begin{figure}[H]
                        \centering
                        \includegraphics[width=0.5\textwidth]{Solution/AMC_8_AMC_2016_25_Solution_Figure_6.png}
                    \end{figure}
                    

 
Letting $17$ be the base of the triangle makes our height the radius of the semicircle:
                    \begin{figure}[H]
                        \centering
                        \includegraphics[width=0.5\textwidth]{Solution/AMC_8_AMC_2016_25_Solution_Figure_7.png}
                    \end{figure}
                    

 Let $r$ be the radius of the semicircle. Then we have $\dfrac{17 \cdot r}{2}=60 \implies 17 \cdot r = 120 \implies \dfrac{17r}{17} = \dfrac{120}{17} \implies r = \boxed{\textbf{(B) }\dfrac{120}{17}}.$

 $\textbf{Remember: Area can help in places where you least expect it!}$

 
 \[\definecolor{salmon}{rgb}{.918, .6, .6}\colorbox{salmon}{\color{white}{solution by \textbf{pog}}}\]

 



\textbf{Solution 8: equations}

Draw tangent lines to connect the centre of the circle to the equal sides. Since the angle is 90 degrees and the height is perpendicular to the side not equal with any of the other side lengths, the two triangles(Not the big one but the two smaller ones) are similar. We've got one common side that gives us the ratio is $\frac{8}{15}$. We've got a system of equations!

 1. 8/15=x(the radius)/y(the side that is similar)2. $x^2 + y^2 = 225$ (by the pythagorean theorem).

 Solving the equations we get $\boxed{\frac{120}{17}}$

 
 \[\definecolor{salmon}{rgb}{.789, .6, .6}\colorbox{salmon}{\color{black}{created by \textbf{justin6688}}}\]

 



\textbf{Solution 9: Process of Elimination}

Not a true solution, but in the interest of creativity and alternative methods, here is a way to cheese the solution in seconds.

 
                    \begin{figure}[H]
                        \centering
                        \includegraphics[width=0.5\textwidth]{Solution/AMC_8_AMC_2016_25_Solution_Figure_2.png}
                    \end{figure}
                    

 First, notice that $AC=17$ (Pythagorean triple $8$-$15$-$17$). Since all the side lengths of $\triangle ADC$ are integers, and $\triangle AED \sim \triangle ADC \sim \triangle DEC$, all of the side length ratios are rational, and therefore, $DE$ must be a rational value. We eliminate answer choices A, D, and E.

 Finally, notice that answer choice C is impossible because that would mean that in $\triangle AED$, the leg $DE = 10$ would be bigger than the hypotenuse $AD = 8$. The only answer choice remaining is $\boxed{\textbf{(B) }\dfrac{120}{17}}$.

 


 


\end{enumerate}